\chapter{Funktionsverbgefüge (FVG)}
\label{cha:FVG}

Funktionsverbgefüge sind Verbindungen aus einem Nomen und einem „leeren" Verb (haben, sein, bringen, stellen), die ein einfaches Verb ersetzen. Sie sind charakteristisch für den formellen, akademischen Stil.

\section{Häufige FVG}
\label{sec:FVGHaeufig}

\begin{center}
\begin{tabular}{|l|l|}
\hline
\textbf{einfaches Verb} & \textbf{FVG} \\ \hline
beeinflussen & Einfluss nehmen auf \\ \hline
kritisieren & Kritik üben an \\ \hline
sich entscheiden & eine Entscheidung treffen \\ \hline
helfen & Hilfe leisten \\ \hline
fragen & eine Frage stellen \\ \hline
beginnen & den Anfang machen mit \\ \hline
enden & ein Ende setzen \\ \hline
diskutieren & zur Diskussion stehen \\ \hline
erlauben & die Erlaubnis geben \\ \hline
beabsichtigen & die Absicht haben \\ \hline
sich bewerben & sich bewerben um \\ \hline
versuchen & den Versuch machen \\ \hline
beweisen & den Beweis erbringen \\ \hline
erklären & eine Erklärung abgeben \\ \hline
verstehen & Verständnis haben für \\ \hline
\end{tabular}
\end{center}

\section{Übungen zu FVG}
\label{sec:FVGUebungen}

\begin{enumerate}
    \item Wir diskutieren die Pläne. $\rightarrow$ \textit{(Lösung: Die Pläne stehen zur Diskussion.)}
    \item Der Arzt kritisiert das Verhalten. $\rightarrow$ \textit{(Lösung: Der Arzt übt Kritik am Verhalten.)}
    \item Sie beginnen mit der Präsentation. $\rightarrow$ \textit{(Lösung: Sie machen den Anfang mit der Präsentation.)}
    \item Er beweist seine Theorie. $\rightarrow$ \textit{(Lösung: Er erbringt den Beweis für seine Theorie.)}
    \item Wir treffen eine Entscheidung. $\rightarrow$ \textit{(Lösung: Wir treffen die Entscheidung.)}
\end{enumerate}

\section{Stilistische Hinweise}
\label{sec:FVGStil}

\begin{tcolorbox}[title=Stil]
\begin{itemize}
    \item \textbf{FVG}: formell, akademisch, schriftlich
    \item \textbf{einfaches Verb}: informell, mündlich
\end{itemize}
\end{tcolorbox}

FVG sollten im akademischen Schreiben verwendet werden, in der gesprochenen Sprache klingen sie jedoch oft steif.

\chapterend